% =============================================================================
% Section VI: Electromagnetic Field
% paper/em_extension.tex
% =============================================================================

\section{Electromagnetic Field}
\label{sec:em}

The decoherence calculations of \cref{sec:powerlaw,sec:desitter} were carried
out for a conformally coupled massless scalar field.  The physically more
important case of an electrically charged body emitting photons is governed
instead by the Maxwell field.  Fortunately, the Maxwell action is conformally
invariant in four spacetime dimensions, so the electromagnetic case follows
from the scalar calculation by a short additional argument.  The main result
of this section is that the electromagnetic decoherence number is exactly
twice the scalar result:
%
\be
  \Nnum_{\rm EM} = 2\,\Nnum_{\rm scalar}
  \label{eq:N-EM-factor2}
\ee
%
for all conformally flat cosmologies, including every member of the power-law
family and de Sitter.  The factor of 2 counts the two independent physical
polarizations of the photon.

\subsection{Conformal invariance of the Maxwell field}
\label{sec:em-conformal}

The Maxwell action in a curved spacetime $(M, g_{ab})$ is
%
\be
  S_{\rm EM}
  =
  -\frac{1}{4}\int F_{ab}\,F^{ab}\,\sqrt{-g}\;\mathrm{d}^4x \,,
  \label{eq:Maxwell-action}
\ee
%
where $F_{ab} = \nabla_a A_b - \nabla_b A_a$ is the electromagnetic field
strength tensor.  Under a Weyl rescaling $g_{ab} \to \Omega^2(x)\,g_{ab}$,
the combination $F_{ab}\,F^{ab}\,\sqrt{-g}$ is invariant in four dimensions:
$F^{ab} = g^{ac}g^{bd}F_{cd}$ acquires a factor $\Omega^{-4}$ while
$\sqrt{-g}$ acquires $\Omega^4$, leaving the product unchanged.  Therefore
\cref{eq:Maxwell-action} is conformally invariant in $d = 4$, and the
equations of motion for $A_a$ are identical in the physical FLRW metric and
in the conformally related Minkowski metric.

In the Lorenz gauge $\nabla^a A_a = 0$, the field equation reduces to
%
\be
  \Box_g A_a = j_a \,,
  \label{eq:Maxwell-eom}
\ee
%
where $j_a$ is the charge current and $\Box_g = g^{ab}\nabla_a\nabla_b$.
Under the conformal rescaling $g_{ab} = \Scfac^2(\eta)\,\eta_{ab}$, the
rescaled potential $\tilde{A}_a \equiv A_a$ (with lower index) satisfies
the flat-space wave equation $\Box_\eta \tilde{A}_a = \tilde{j}_a$ — the
conformal weight of the Maxwell field is zero for a covariant vector in four
dimensions~\cite{Wald84}.  The free in-modes of $A_a^{\rm in}$ are thus
in one-to-one correspondence with the flat-space photon modes, and the
conformal vacuum for the Maxwell field is again the pullback of the Minkowski
photon vacuum.

\subsection{Electromagnetic Wightman function and decoherence}
\label{sec:em-deco}

The electromagnetic Wightman 2-point function in the conformal vacuum is
%
\be
  \langle\Psi_0|\,A_a^{\rm in}(x_1)\,A_b^{\rm in}(x_2)\,|\Psi_0\rangle
  =
  W_{ab}^{\rm Mink}(x_1, x_2) \,,
  \label{eq:EM-Wight}
\ee
%
where $W_{ab}^{\rm Mink}$ is the flat-space photon propagator in Lorenz gauge.
The FLRW conformal factors drop out because the covariant potential $A_a$ has
conformal weight zero, in contrast to the scalar field $\varphi$ of
\cref{eq:Wight-conformal-relation} whose conformal weight is $-1$.  (The
scalar $\varphi$ and the photon $A_a$ carry different conformal weights, but
both fields are conformally invariant in the sense that their dynamics in FLRW
are governed by the flat-space wave operator.)

For a charge $q$ whose center-of-mass is in a superposition of worldlines
$\gamma_1$ and $\gamma_2$, the electromagnetic current difference is
%
\be
  j_1^a(x) - j_2^a(x)
  \approx
  q\, d(\tau)\, s^a\, \nabla_b\, \delta^{(3)}\!\lb(\mb{x} - \mb{X}(\tau)\rb) \,,
  \label{eq:EM-current-dipole}
\ee
%
in the same dipole approximation as \cref{eq:dipole}.  Substituting
\cref{eq:EM-current-dipole} into the electromagnetic analogue of
\cref{eq:Gamma-exact},
%
\be
  \Gamma_{\rm EM}
  =
  \langle\Psi_0|\,
    \lb[A_a^{\rm in}(j_1^a - j_2^a)\rb]^2
  \,|\Psi_0\rangle \,,
  \label{eq:Gamma-EM}
\ee
%
we find that the integrand factorizes into a transverse-polarization sum and
the scalar Wightman kernel.  Explicitly, the flat-space photon propagator in
Lorenz gauge at coincident spatial points takes the form
%
\be
  W^{\rm Mink}_{ij}(\eta, \mb{X};\, \eta', \mb{X})
  =
  \Pi_{ij}^{\rm T} \cdot W_{\rm Mink}(\eta, \mb{X};\, \eta', \mb{X}) \,,
  \label{eq:EM-polarization-sum}
\ee
%
where $\Pi_{ij}^{\rm T} = \delta_{ij} - \hat{x}_i \hat{x}_j$ is the
transverse projector averaged over the direction of emission, and
$W_{\rm Mink}$ is the scalar Wightman function of \cref{eq:Wight-conformal-relation}.
Contracting with the dipole source $s^i s^j$ and summing over the transverse
directions,
%
\be
  s^i s^j\,\Pi_{ij}^{\rm T}
  =
  s^i s^j \lb(\delta_{ij} - \hat{k}_i \hat{k}_j\rb)
  \label{eq:polarization-contraction}
\ee
%
and integrating over all emission directions $\hat{k}$, one obtains
$\int \mathrm{d}^2\hat{k}\, s^i s^j(\delta_{ij} - \hat{k}_i\hat{k}_j) = 2\times$
the corresponding scalar integral~\cite{DSW2022,Danielson:2022tdw}.  The factor
of 2 counts the two independent transverse polarization states of the photon:
for each propagation direction, a massless spin-1 field has two physical
degrees of freedom, compared to the single degree of freedom of the scalar.
Therefore,
%
\be
  \Gamma_{\rm EM} = 2\,\Gamma_{\rm scalar} \,,
  \label{eq:Gamma-EM-2}
\ee
%
and consequently $\Nnum_{\rm EM} = 2\,\Nnum_{\rm scalar}$, establishing
\cref{eq:N-EM-factor2}.

\subsection{Summary of electromagnetic results}
\label{sec:em-summary}

Combining \cref{eq:N-EM-factor2} with the scalar results of
\cref{sec:powerlaw,sec:desitter}, the electromagnetic decoherence numbers are:

\begin{itemize}[noitemsep]
  \item \emph{De Sitter spacetime:}
  %
  \be
    \Nnum_{\rm EM}^{\rm dS}
    =
    \frac{q^2 d_0^2 H^3}{48\pi^2}\, T \,,
    \label{eq:N-EM-dS}
  \ee
  %
  linear in the superposition time $T$, with the thermal prefactor
  proportional to $T_{\rm GH}^3 = (H/2\pi)^3$.

  \item \emph{Power-law FLRW with $p > 1$:}
  %
  \be
    \Nnum_{\rm EM}^{\rm power-law}
    \;\sim\;
    \frac{q^2 d_0^2 p^3}{24\pi^3\,T_0^2}\,\ln\!\lb(\frac{T}{T_0}\rb) ,
    \qquad T \gg T_0 \,,
    \label{eq:N-EM-PL}
  \ee
  %
  logarithmic for large $T$, with the same non-thermal prefactor as the
  scalar result and no Gibbons-Hawking temperature structure.
\end{itemize}

The universality of the factor of 2 across all conformally flat FLRW backgrounds
is a direct consequence of conformal invariance: the photon polarization count
is a kinematic property of the field that is independent of the background
geometry.  We collect all decoherence results — scalar and electromagnetic, de
Sitter and power-law — in Table~\ref{tab:summary} of \cref{sec:discussion}.

%% CITATIONS NEEDED (for gpd-bibliographer):
%% - Wald84        -- Wald, General Relativity (1984) [conformal weights]
%% - DSW2022       -- Danielson, Satishchandran, Wald (2022b) IJMPD 31, 2241003
%% - Danielson:2022tdw -- DSW (2023) PRD 108, 025007
